% Generated draft: requires mathematical and source review.
\documentclass[11pt]{article}
\usepackage[margin=25mm]{geometry}
\usepackage{fontspec}
\setmainfont{Latin Modern Roman}
\usepackage{amsmath,amssymb,mathtools,unicode-math}
\setmathfont{Latin Modern Math}
\usepackage{enumitem,needspace,xurl,hyperref}
\setlength{\parindent}{0pt}
\setlength{\parskip}{5pt}
\setlength{\emergencystretch}{3em}
\newcommand{\N}{\mathbb{N}}
\newcommand{\Z}{\mathbb{Z}}
\newcommand{\Q}{\mathbb{Q}}
\newcommand{\R}{\mathbb{R}}
\newcommand{\C}{\mathbb{C}}
\newcommand{\F}{\mathbb{F}}
\newcommand{\A}{\mathbb{A}}
\newcommand{\PP}{\mathbb{P}}
\newcommand{\E}{\mathbb{E}}
\newcommand{\eps}{\varepsilon}
\newcommand{\dd}{\,\mathrm{d}}
\begin{document}
\section*{UnsolvedMath ID: 30006013}
\textbf{Source problem code:} catalogue entry
\subsection{Definitions and mathematical statement}
The supplied record poses two related research questions, but does not supply the definitions and hypotheses needed for a fully specified mathematical statement.

\begin{enumerate}
\item Characterize the algebraic stacks satisfying the \emph{cyclic lifting property}. The record does not define this property, its lifting diagrams, the cyclic groups involved, or the ambient class of stacks and base schemes. Consequently, its quantifiers and necessary and sufficient conditions cannot be stated precisely from the supplied evidence.
\item Formulate and establish an extension of the associated \(p\)-adic volume formula, described in the source as an analogue of Weil's formula, for moduli spaces having non-quotient singularities. The examples named are moduli spaces of Higgs bundles, coherent sheaves on projective varieties, and quiver representations. The record supplies neither the original formula nor the proposed extended identity, integration measure, integrand, convergence conditions, or full hypotheses on these moduli spaces.
\end{enumerate}

The source suggests a starting setup involving a linear quotient stack
\[
\mathcal M=[U/\mathrm{GL}_N],
\]
where \(U\) is a smooth algebraic \(\mathcal O_F\)-space. Here the quotient notation requires an action of \(\mathrm{GL}_N\) on \(U\); the action is not specified in the record. The source also describes a map that is an isomorphism over a dense open subset whose complement has codimension at least two. If the two occurrences of \(M\) in the extracted map are intended to distinguish the quotient stack from a target space, this condition can be written as
\[
\pi:\mathcal M\longrightarrow M,\qquad
W\subseteq M\text{ dense and open},\qquad
\mathcal M\times_M W\xrightarrow{\sim}W,\qquad
\operatorname{codim}_M(M\setminus W)\geq 2.
\]
This is a conditional transcription: the extracted source does not reliably distinguish the stack and target notation, identify the nature of \(M\), or specify \(F\), \(\mathcal O_F\), or the allowable ranks \(N\). These conditions are presented as a first step, not as sufficient hypotheses for an asserted formula.

The source reports an application in its reference [6] to a moduli Deligne--Mumford stack of \(G\)-Higgs bundles, subject to additional assumptions ensuring smoothness. Those assumptions are absent. It further reports a description of the twisted inertia stack as a disjoint union indexed by endoscopic groups of \({}^{L}G\); the notation for the constituent moduli spaces is corrupted in the supplied extraction and is not reconstructed here. The source connects the \(p\)-adic volume of the coarse moduli space and related integrals with a proof of the fundamental lemma, referring to Ng\^o's reference [8]. These are supplied background claims, not additional theorems asserted by this draft.

For the extension to non-quotient singularities, the source describes ongoing work and reports a satisfactory understanding for moduli spaces of objects in hereditary abelian categories with symmetric Euler forms, subject to further assumptions explicitly omitted there. This qualification must be retained: the record does not support an unconditional assertion for all such categories or all the listed moduli spaces.
\subsection{Short English statement}
Determine which algebraic stacks have the cyclic lifting property, and extend the associated \(p\)-adic volume formula to suitable moduli spaces with non-quotient singularities. The supplied record omits the property's definition, the formula, and essential hypotheses, so this remains a qualified research question requiring human mathematical review.
\subsection{Sources}
\begin{itemize}
\item Supplied local UnsolvedMath registry record, permanent ID \(30006013\), title \emph{Cyclic Lifting Property for Moduli Stacks}, exact source problem code \texttt{catalogue entry}; collection: \emph{Oberwolfach Reports Open Problems}. No individual proposer is identified.
\item Supplied source citation: \emph{Arithmetic Geometry (2025)}, Oberwolfach Reports, DOI \texttt{10.4171/owr/2024/33}, pp.\ 1855--1912. Supplied source URL and literature-source link, labelled ``OWR report'': \url{https://doi.org/10.4171/owr/2024/33}. The record separately gives proposed year 2024; the discrepancy with the citation's title year is unresolved.
\item Literature assessment dated exactly \texttt{2026-08-22}: the supplied curation labels the problem open and states that the cyclic lifting property is verified in the examples used for the volume formula, while its general range of validity is left open. This is supplied dated curation, not an independently verified literature assessment or source verification.
\item The extracted passage cites its reference [6] for the application to \(G\)-Higgs bundles and Ng\^o's reference [8] for the fundamental lemma. Their bibliographic entries are not supplied; no titles, authors for [6], or additional URLs are inferred.
\end{itemize}
\end{document}
